Scientific discovery often requires optimising costly-to-evaluate objectives over vast, structured, and open-ended hypothesis spaces, such as molecules, protein sequences, and computer programs. Large language models (LLMs), or domain-specific generative models, provide flexible generative priors over these spaces. However, their likelihoods and self-assessments do not provide reliable estimates of empirically measured results or calibrated epistemic uncertainty, particularly for novel candidates outside the observed data distribution. We propose the Large Discovery Model (LDM), an empirically grounded recurrent architecture in which an LLM generates and refines candidate designs, while a continually updated Bayesian non-parametric reward model, instantiated here as a Gaussian process (GP), predicts their quality and quantifies uncertainty. A GP-derived acquisition function guides candidate generation, refinement, and selection, while new observations update the surrogate and discovery memory. We formulate this process as a KL-regularised, acquisition-guided search objective and implement it through candidate reweighting, iterative refinement, and decomposed feedback on predicted quality, uncertainty, and constraint satisfaction. We evaluate LDM on neural-network training-program search, antibody CDRH3 sequence design, and multi-objective molecular optimisation. These studies demonstrate how LDM integrates structured generative capabilities of LLMs with uncertainty-aware, experiment-grounded search to support open-ended discovery across diverse scientific design spaces.

%Scientific discovery often requires optimising costly-to-evaluate objectives over vast, structured, and potentially open-ended spaces, such as molecules, protein sequences, and computer programs. Large language models (LLMs), or other domain-specific generative models, provide flexible generative priors over these spaces but cannot reliably estimate objective values or uncertainty. We propose the Large Discovery Model (LDM), an experiment-grounded recurrent architecture in which an LLM generates and refines candidate designs, while a continually updated Bayesian non-parametric reward model, instantiated here as a Gaussian process (GP), predicts their quality and uncertainty. A GP-derived acquisition function guides candidate generation, refinement, and selection, while new observations update the surrogate and discovery memory. We formulate this process as a KL-regularised, acquisition-guided search objective and implement it through candidate reweighting, iterative refinement, and decomposed feedback on predicted quality, uncertainty, and constraint satisfaction. We evaluate LDM on neural-network training-program search, antibody CDRH3 sequence design, and multi-objective molecular optimisation. These different studies demonstrate how LDM achieves the structured generative capabilities of LLMs with uncertainty-aware, experiment-grounded search to support open-ended discovery across diverse scientific design spaces.

%\vspace{0.5em}
This initial release (LDM v0.1) accompanies the public framework, code, and benchmarks available at \url{https://github.com/yzailab/Large-Discovery-Models}. 
%Future versions of this report will document new capabilities, experiments, and updates to the framework.

